\section{相关工作与研究背景}
\subsection{临床长文本建模}
临床文本任务与通用文本分类任务相比，具有更强的长程依赖、更多的领域术语以及更复杂的证据分布。单条病历中往往同时包含主诉、现病史、既往史、检验检查、医嘱和随访等多个部分，模型若仅依赖局部片段，容易遗漏跨段落、跨时间形成的关键信息。因此，能够稳定处理长上下文的医学编码器为临床自然语言处理提供了重要基础。本文采用的 \encoder{} 正是建立在这一背景之上的长上下文编码器，其作用是为后续预测模块提供质量较高的全局语义表示 \citep{sounack2025bioclinicalmodernbertstateoftheartlongcontext}。

然而，长上下文编码能力并不自动等同于高质量决策能力。对于临床任务而言，模型不仅需要“看到足够多的信息”，还需要在证据不完整、标签边界模糊和样本异质性明显的条件下做出稳健判断。基于这一认识，本文不把研究重点继续放在编码器主干的扩展上，而是进一步讨论在强表示已经可得的前提下，输出层还能承担哪些建模职责。

\subsection{临床预测中的不确定性问题}
临床预测任务普遍伴随显著的不确定性来源。其一，临床记录本身具有书写噪声，不同医生、科室和时间段的记录风格存在差异，导致同一标签对应的文本证据形态并不统一。其二，许多临床标签具有弱监督或边界模糊特征，例如共病状态、风险倾向和行为表型等标签，往往需要综合多处间接线索才能判断。其三，类别不平衡和难例稀缺广泛存在，模型在平均意义上表现良好，并不意味着其在高风险样本上同样可靠。

因此，在临床场景中，仅报告单一准确率或 F1 指标通常是不充分的。模型若无法区分“证据充分的容易样本”和“证据不足的高风险样本”，就难以在真实流程中支持阈值控制、人工复核和风险分级。不确定性建模的意义，正在于为预测结果附加一个能够反映样本难度或噪声水平的辅助信号，从而使模型输出更接近临床使用中的决策需求。

\subsection{深度表示与结构化输出层}
在许多深度学习应用中，研究者往往把主要精力放在编码器或主干网络上，而把输出层视为相对简单的线性映射或浅层 MLP。这种做法在标准数据集上通常足够有效，但在临床文本场景中未必总是充分。一方面，输出层直接决定了深度特征如何转化为最终预测，其归纳偏置会影响模型对噪声样本和边界样本的处理方式；另一方面，输出层也是最容易与传统统计估计方法结合的位置，适合引入权重、正则化和解析求解等结构信息。

从这一角度看，结构化输出层并不是对深度模型的简单补丁，而是连接“高维表示学习”和“稳健统计估计”的关键界面。若能够在该界面上同时建模预测值与样本相关噪声，并利用噪声信息重估最终映射关系，就有可能在不改动主干编码器的前提下改善模型的稳健性、解释性以及跨样本的一致性。

\subsection{本文工作的定位}
基于上述背景，本文将研究问题界定为：在冻结临床长文本编码器的条件下，能否通过面向输出层的不确定性建模与结构化精修，获得比常规线性头和 MLP 头更稳定的综合表现。为回答这一问题，本文选择在特征空间中联合建模预测均值与输入相关方差，并进一步利用方差诱导的样本权重对最后一层参数进行加权最小二乘精修。这样的设计既保持了深度特征的表达能力，又引入了可分析的统计结构，因此与单纯增加隐藏层深度的做法形成了明确区别。

本文工作的意义不在于重新提出一个更大的基础模型，而在于说明：当强编码器已经存在时，输出层仍然是值得独立研究的对象。尤其在临床文本这类高噪声、高异质任务中，对输出层进行可信建模，不仅可能改善指标，也有助于为后续校准、拒识和人工协同提供接口。
